\section{Introduction}
\label{sec:introduction}

Binary encoding is a major computation in correlation generators and
code-based proof systems. These applications need constant rate and linear
distance, but their cost also depends on how the encoder moves data and
how efficiently its transpose can be evaluated.

Pseudorandom correlation generators apply a public generator matrix to large
vectors to expand sparse correlations into OT and VOLE correlations
\cite{bcgi18,bcg19pcg}. Code-based proof systems encode rows of large arrays,
with distance supporting soundness~\cite{brakedown23,blaze25}.
At these dimensions, long-range memory traffic can cost
more than the binary arithmetic.

Random linear codes provide excellent distance, but direct encoding requires
dense computation. Structured encoders can reduce both arithmetic and memory
traffic. Proving strong distance for such an encoder requires controlling
the dependencies introduced by its structure. Our goal is a binary code
family that combines a quantitative distance guarantee with linear-time
evaluation in both directions and regular data movement.

We introduce Single-Permutation INterleaved (SPIN) codes. Our main result,
Structured SPIN, meets this goal. At its native output lengths
$N$, every realized code has rate $1/2$ and $O(N)$ ordinary and transposed
encoding work. With probability $1-o(1)$ over setup, its minimum distance
exceeds $0.11N$. The construction combines a proved distance guarantee
with an encoding structure designed for efficient implementation. We
evaluate this combination both as an encoder for PCGs and as a component
of code-based proof systems.

SPIN follows the serial-concatenation pattern. An outer encoder first
prevents nonzero messages with small support from remaining too sparse.
The interleaver then scatters the outer support across the input to the
inner encoder. This reduces the chance that local structure in the outer
word causes the inner encoder to produce a low-weight output.
Finally, the inner encoder generates the codeword by recursive convolution.
SPIN makes this pattern
modular: structure appears where speed and regular data movement require it,
and setup randomness remains at the interfaces used by the distance proof.

For an admissible output length $N$, a SPIN encoder has the form
\begin{equation}
  E_N := I_N\circ\Pi_N\circ O_N.
  \label{eq:intro-spin-composition}
\end{equation}
The outer encoder $O_N$ maps the message into $N$ intermediate bits.
The interleaver $\Pi_N$ is one permutation of those $N$ coordinates.
The inner encoder $I_N$ is a length-preserving recursive linear map.
Setup may sample any of these objects from a declared distribution.

The phrase \emph{single permutation} describes the mathematical interface.
An implementation may factor $\Pi_N$ into block shuffles, a transpose, and
region shuffles. These factors still compose to one global interleaver.

Beyond the code construction, we evaluate SPIN in two applications:
correlation generation and code-based proof systems. For PCGs, we measure
the encoding stage and estimate the resulting OT throughput. For proof
systems, we instantiate the Brakedown matrix polynomial commitment scheme
(PCS)~\cite{brakedown23} with SPIN, then integrate it into
Flock~\cite{flock26}. Standalone commitment and opening benchmarks measure
the PCS cost; the Flock comparison measures its effect on a complete
constraint prover, using Ligerito~\cite{ligerito25} as the alternative PCS.
The resulting SPIN-based PCS favors prover speed over proof size.

\subsection{Three SPIN constructions}

The three constructions develop the same idea in stages: start with a
simple accumulator, strengthen the inner mixing, then replace expensive
random components with efficient structured ones.

\paragraph{Accumulator SPIN.}
We begin with random outer blocks, one uniform interleaver, and a binary
accumulator. This simplifies the Block-Accumulate (BA) architecture,
including the Block-Accumulate-Accumulate (BAA) constructions of prior
work~\cite{KolesnikovEtAl2026BA,cryptoeprint:2026/1903}, to one
interleaver and one accumulator. The accumulator has an exact input--output
weight enumerator, making it a useful starting point for the analysis.
However, a single accumulator provides limited mixing of sparse inputs.
Our rate-$1/2$ theorem gives relative distance $0.0037$, far below the
Gilbert--Varshamov (GV) benchmark of approximately $0.110028$.
This weak guarantee motivates a stronger inner encoder.

\paragraph{Random SPIN.}
We retain the random outer blocks and uniform interleaver, but replace
the accumulator with a random recursive convolution of logarithmic memory.
The larger memory lets the inner mix sparse inputs more effectively.
With logarithmic outer block size, our rate-$1/2$ theorem proves relative
distance greater than $0.11002$, nearly matching the GV bound, with
$O(N\log N)$ direct encoding work.

This recursive inner has a lineage in codes for PCGs. Silver used a banded
lower-triangular solve, but its linear-distance claim remained
conjectural~\cite{CouteauRindalRaghuraman2021Silver}.
Expand--Convolute introduced a random recursive convolution and proved
distance bounds through Markov-chain coupling and concentration
\cite{RaghuramanRindalTanguy2023EC}. The proved parameters are conservative;
its evaluation also proposed faster, aggressive parameters supported by
empirical distance estimates beyond those guarantees. We retain its random inner but analyze the
concatenation through weight enumerators. An exact input--output transform
tracks the expected weight spectrum and yields the near-GV distance bound.

\paragraph{Structured SPIN.}
Random SPIN establishes that the architecture can approach the GV bound,
but its dense random outer blocks and random convolution are expensive.
Structured SPIN identifies the properties needed by that proof and realizes
them with efficient constituent codes and structured interleaving. The
outer supplies bounds on the number of codewords at each weight; the inner
supplies bounds on how input weights propagate to output weights.
These interfaces let us replace the random components while retaining
strong distance guarantees. Our scalable family has rate $1/2$, relative
distance greater than $0.11$ with high probability, and $O(N)$ encoding work.

The same analysis supports concrete parameter selection using exact
enumerators where available and rigorous spectrum bounds otherwise.
For the implemented family, we certify relative distance greater than
$0.10$ with setup failure below $2^{-40}$ at five message lengths from
$2^{16}$ to $2^{24}$. Parameter selection trades encoding cost against the
desired setup-failure bound, with each choice checked by its certificate.
The resulting encoders are several times faster than prior implementations,
as quantified below. This connects the asymptotic distance argument to
finite guarantees and practical encoding speed.

Table~\ref{tab:spin-families} summarizes the three lines.

\begin{table}[t]
  \centering
  \caption{The three principal SPIN construction lines.}
  \label{tab:spin-families}
\begin{tabular}{@{}llll@{}}
    \toprule
    Family & Outer & Interleaver & Inner \\
    \midrule
    Accumulator SPIN & random blocks & uniform & accumulator \\
    Random SPIN & random blocks & uniform & random recursive \\
    Structured SPIN & spectrum-controlled & structured & structured recursive \\
    \bottomrule
  \end{tabular}
  \end{table}

The interface also permits other combinations of outer encoder,
interleaver, and inner encoder. Our analysis treats the three families above.

\subsection{Contributions}

We make five contributions.

\begin{itemize}
  \item \textbf{A structured linear-time code with proved distance.}\par
  We construct a sampled binary code family whose native members have rate
  $1/2$ and minimum distance greater than $0.11N$ with probability $1-o(1)$.
  Every realized member has $O(N)$ ordinary and transposed encoding work. The
  construction combines a spectrum-controlled constituent reused along the
  outer diagonal, one factored global interleaver, and a structured recursive inner.
  Zero extension gives every sufficiently large requested output length
  with an $o(1)$ loss in rate and relative distance.

  \item \textbf{A first-moment framework for structured interleavers.}\par
  We group outer words into classes and bound each class's conditional
  failure probability under the interleaver--inner pair. This accommodates
  structured routing, where equal-weight words need not have the same
  output-weight distribution. Accumulator SPIN and Random SPIN isolate the
  roles of the outer spectrum, interleaving, and recursive mixing. At rate
  $1/2$, they prove relative distances $0.0037$ and $0.11002$, respectively.

  \item \textbf{Finite certificates and measured encoding.}\par
  We combine a fixed BCH-derived $[256,128,\ge38]$ constituent with a
  selected $(128,19)$ RM2Sub inner. We certify relative distance greater
  than $0.10$ with setup failure
  below $2^{-40}$ at message lengths $2^{16},2^{18},2^{20},2^{22},2^{24}$.
  The proof uses rigorous spectrum inequalities despite the unknown exact
  BCH spectrum. Separate paired benchmarks of $128$ parallel instances
  show similar ordinary and transposed encoding costs: about $11$\,ms
  at message length $K=2^{20}$ on one Ryzen 9 7950X thread.

  \item \textbf{A SPIN-based polynomial commitment implementation.}\par
  We instantiate Brakedown with SPIN and organize the computation around
  packed binary data. At $512$\,MiB input, commitment and one opening take
  $520$\,ms on one core. Standalone Ligerito measurements show the tradeoff:
  SPIN--Brakedown proves faster, while Ligerito has smaller proofs and faster
  verification. SPIN--Brakedown's commitment is $4.16\times$ faster than the fastest
  measured Bolt-max configuration in a separate commitment-only experiment
  at the same input volume. Section~\ref{sec:spin-pcs} states the
  conditional parameter target and comparison scope.

  \item \textbf{Integration into Flock.}\par
  We adapt the PCS to Flock's witness claims and optimize the shared prover
  implementation. With the applicable optimizations in both configurations,
  SPIN--Brakedown improves throughput by $1.93\times$ and $1.38\times$
  over Ligerito at the two measured workloads, with larger proofs.

\end{itemize}

\subsection{Related work}

\paragraph{Serially concatenated codes.}
Turbo codes introduced randomized interleaving as a way to combine simple
constituent encoders~\cite{berrou93}. Benedetto et al. developed the
uniform-interleaver analysis of serial concatenations~\cite{benedetto98}.
Divsalar, Jin, and McEliece introduced repeat--accumulate codes and their
weight-enumerator analysis~\cite{divsalar98}, while Kahale and Urbanke studied
minimum distance in random parallel and serial concatenations
~\cite{kahale98}. Bazzi, Mahdian, and Spielman proved limitations for
depth-two constant-memory serial concatenations and positive results with two
accumulators~\cite{bms09}. Later work established coding theorems for
repeat-multiple-accumulate ensembles approaching the binary
Gilbert--Varshamov bound~\cite{kliewer08}.

Block-Accumulate codes replace repetition with random block-diagonal maps and
compute detailed expected spectra; they also provide optimized CPU and GPU
implementations~\cite{KolesnikovEtAl2026BA}. Accumulator SPIN is their
one-accumulator truncation. The newer companion work on chosen-block BAA
updates this family by reusing one spectrum-controlled constituent in every
diagonal position~\cite{cryptoeprint:2026/1903}. Structured SPIN combines
this shared design idea with a structured global route and a different
recursive inner. Our benchmark table compares these newer companion
instantiations; the abstract's speedup is relative to original BAA.

\paragraph{Codes for correlation generation.}
Linear codes compress sparse correlations in vector-OLE and pseudorandom
correlation generators~\cite{bcgi18,bcg19pcg}. This application motivated a
sequence of faster structured codes. Silver uses structured LDPC matrices and
a banded triangular solve~\cite{CouteauRindalRaghuraman2021Silver}. Linear
distance was an explicit conjecture supported by heuristic evidence.
Expand--Accumulate combines a sparse expander with an accumulator
~\cite{ea22}. Expand--Convolute replaces that accumulator with a random
time-varying convolution~\cite{RaghuramanRindalTanguy2023EC}. Random SPIN keeps
the latter inner distribution but replaces the expander with random block
injections. Structured SPIN instead combines a spectrum-controlled outer,
one factored global interleaver, and the RM2Sub inner. Its principal distinction
is the simultaneous proof of rate one half, relative distance $0.11$, and
linear ordinary and transposed encoding work.

\paragraph{Linear-time proving systems.}
Constant-overhead arguments make linear-time encoding a first-order prover
cost~\cite{rr25}. Orion uses code switching to obtain a linear-time prover
~\cite{orion22}, and Brakedown builds a field-agnostic SNARK from a
linear-time encodable code~\cite{brakedown23}. Field-agnostic SNARKs have also
been instantiated from Expand--Accumulate codes~\cite{fieldEA24}. Blaze uses
interleaved RAA codes and supplies improved binary distance bounds
~\cite{blaze25}. Its generic code-switching framework also accommodates
SPIN in theory, with asymptotically smaller proofs than our measured
Brakedown-style construction. Bolt uses sketched codes to reduce commitment
cost~\cite{bolt26}. We instantiate the Brakedown matrix PCS with SPIN
and evaluate it independently before integrating it into Flock's Boolean
constraint prover~\cite{flock26}. The latter comparison uses Ligerito
as the alternative PCS~\cite{ligerito25}.

\subsection{One distance calculation}

All three lines use the same first-moment idea. Fix a target distance $d$,
and count nonzero messages whose encoded words have weight at most $d$.
If the expected count is below one, then some sampled encoder is injective and
has minimum distance greater than $d$.

For a uniform interleaver, the outer Hamming weight is often a sufficient
interface. The expected bad-word count then has the form
\begin{equation}
  \sum_w A_w^{\mathrm{out}}p_w^{\mathrm{in}}(d),
  \label{eq:intro-weight-first-moment}
\end{equation}
where $A_w^{\mathrm{out}}$ is the expected outer spectrum and
$p_w^{\mathrm{in}}(d)$ is an inner slice-to-tail probability.

A structured interleaver need not produce a uniform Hamming slice. We
therefore allow a class to record block occupation or another routing profile.
The interleaver--inner analysis supplies a uniform failure envelope for each
class. The weight-indexed sum in~\eqref{eq:intro-weight-first-moment} is the
special case in which the conditional law depends only on weight.

\subsection{Organization}

Section~\ref{sec:preliminaries} fixes notation. Section~\ref{sec:framework}
states the exact finite first-moment framework.
Sections~\ref{sec:accumulator-spin} and~\ref{sec:random-spin} develop the two
random construction lines. Section~\ref{sec:structured-spin} defines the
  structured line, its parameterized constituent interface, and the currently
  certified scalable instantiation.
Section~\ref{sec:finite-certificates} presents the finite BCH construction,
its distance certificates, and the certified engineering curve.
Section~\ref{sec:scaling} separates finite validity, asymptotic distance, and
encoding complexity. Section~\ref{sec:implementation} records the
encoder implementation and measurements. Section~\ref{sec:spin-pcs}
presents the standalone PCS, and Section~\ref{sec:spin-flock} evaluates
its integration into Flock. The appendices give the asymptotic and
finite proof reductions and identify their computational evidence.

\paragraph{Artifact.}
\ifsubmission
The source repository includes an artifact guide and a paper-to-code map.
\else
The source repository is
\artifactlink{https://github.com/ladnir/permute_conv}{\texttt{ladnir/permute\_conv} on GitHub}.
Its \artifactlink{https://github.com/ladnir/permute_conv/blob/main/artifact/README.md}{artifact guide}
gives the reproduction entry point, and its
\artifactlink{https://github.com/ladnir/permute_conv/blob/main/artifact/PAPER_MAP.md}{paper-to-code map}
identifies the corresponding sources.
\fi
The map links the theorems, figures, and timing tables to their sources
and checks. The guide separates reproduction of reported values from
numerical proof replay and new performance measurements.
